\chapter{Overlap neighborhoods and reference configurations}\label{sec:references}

This section collects the probabilistic input of the dynamical argument in
\cref{sec:induction,sec:barrier-proof}. It has three parts. In
\cref{sec:ref-neighborhoods}, we prove \cref{b:lem:coverage-extension}, the
overlap-neighborhood mass bound. It states that outside an event of
vanishing disorder probability, for every set $S\subseteq\Sigma_n$, every
$x\in S$ outside a set of Gibbs mass less than $e^{-D_1L_n}$ has the property
that the configurations $z\in S$ with $|R(x,z)|\ge q_*$ carry Gibbs mass at
least $e^{-D_3L_n}$;
here $D_1>0$ is arbitrary, and $D_3$ depends only on $\beta$ and $D_1$.
This is the argument of \cite[Proposition~6.1]{oai-barriers}, carried out at
every diverging sublinear scale $L_n$. In \cref{sec:ref-parameters}, we fix the
parameters of the dynamical argument and the order in which they are chosen
(\cref{tab:ref-parameters}). We then define the set $\mathcal B_n$ of
exceptional triples as an explicit union and deduce from
\cref{b:lem:static-event} that $\pi^{\otimes3}(\mathcal B_n)$ is small on
one disorder event (\cref{lem:ref-single-event}). In
\cref{sec:ref-families}, we combine these two inputs. At each disorder in an
event of probability tending to one, we fix one family of Gibbs reference
configurations and one set $V$ of small Gibbs mass such that the four
properties used in \cref{sec:induction} hold along every path that avoids
$V$ (\cref{lem:ref-families}).

The section uses \cref{b:lem:static-event} and the pressure estimate
\eqref{b:eq:fixed-pressure-gap} from \cref{sec:static}, the scalar
constant of \eqref{b:eq:main-analytic-constants} from \cref{sec:scalar},
the band of \eqref{b:eq:os-band-checks} from \cref{sec:cubic}, and the
imported results \cref{inp:ob-energy-tail,inp:ob-rotation,inp:ob-concentration}.
\Cref{lem:ref-families} is the only point at which the static
estimates of \cref{part:static} enter the induction of \cref{sec:induction};
its bound on $\pi(V)$ enters the stationary sign test of
\cref{sec:barrier-proof}. Although \cref{thm:barrier-730} uses only the scale
exponent $7/30$, we state the construction for general exponents, because this
shows which conditions on the scales each step uses
(see \cref{tab:bp-exponents}).

Throughout this section, $\beta>1$ is fixed, and we write $\pi=\pi^J$. For
a vector $\bar J=(\bar J_{ij})_{1\le i<j\le n}$ of real numbers, we write
$H_n^{\bar J}$ and $Z_n^{\bar J}$ for the Hamiltonian and the partition
function in which $J$ is replaced by $\bar J$. For $k\ge1$ and a
set $\mathcal U\subseteq\Sigma_n^k$, $\pi^{\otimes k}(\mathcal U)$ is the
probability, at fixed disorder, that $k$ independent samples from $\pi$ form
a $k$-tuple in $\mathcal U$; we also write
$\pi^{\otimes2}(|R|\ge q)$ for
$\pi^{\otimes2}(\{(x,y):|R(x,y)|\ge q\})$. The Gibbs measure is invariant
under the spin flip $x\mapsto-x$, because $H_n^J(-x)=H_n^J(x)$, and $\pi^{\otimes k}$
is invariant under spin flips of single coordinates and under permutations of
the coordinates. All constants depend on $\beta$ and on the fixed parameters
listed in \cref{tab:ref-parameters}, and never on $n$ or on the disorder.

\section{Overlap neighborhoods at sublinear scales}\label{sec:ref-neighborhoods}

The bound below was proved in \cite[Proposition~6.1]{oai-barriers} at the
scale $L_n=n^{1/10000}$ used there; this is \cref{inp:ob-coverage}. Its proof
uses the scale only through the conditions $L_n\to\infty$ and $L_n=o(n)$. Our
time scale is different, and the constants of the bound enter the
construction of \cref{sec:ref-families}. For these reasons we state the
extension and write out the proof. The ingredients are those of
\cite{oai-barriers}: the upper
tail of the energy of a Gibbs sample (\cref{inp:ob-energy-tail}, whose proof
in \cite[Lemma~5.1]{oai-barriers} uses the Brascamp--Lieb inequality
\cite{brascamp-lieb-1976} in the form of \cite[Lemma~2.1]{oai-barriers}), the
disorder rotation estimate (\cref{inp:ob-rotation}), Gaussian concentration
for Lipschitz functions (\cref{inp:ob-concentration}), and the
Paley--Zygmund inequality.

\begin{lemma}[Overlap-neighborhood mass bound at sublinear scales]
\label{b:lem:coverage-extension}
There exists $q_*=q_*(\beta)\in(0,1)$ with the following property. Let
$(L_n)$ be a deterministic sequence of positive numbers with $L_n\to\infty$
and $L_n/n\to0$, and let $D_1>0$. Then there exist $D_3>0$ and $D_4>0$, depending only
on $\beta$ and $D_1$, and disorder events $\mathcal C_n(D_1)$ such that
$\PP(\mathcal C_n(D_1)^c)\le3e^{-D_4L_n}$ for all sufficiently large $n$,
and such that on $\mathcal C_n(D_1)$,
\begin{equation}\label{b:eq:general-coverage}
 \pi\bigl(\bigl\{x\in S:
   \pi(\{z\in S:|R(x,z)|\ge q_*\})<e^{-D_3L_n}\bigr\}\bigr)
 <e^{-D_1L_n}
 \qquad\text{for every }S\subseteq\Sigma_n .
\end{equation}
\end{lemma}

Both masses in \eqref{b:eq:general-coverage} are masses under $\pi$ itself,
not under $\pi$ conditioned on $S$. The inner mass is the probability that a
single Gibbs sample lands in the overlap neighborhood
$\{z\in S:|R(x,z)|\ge q_*\}$ of $x$ in $S$; in \cref{sec:ref-families} it is
the success probability of one reference sample. The overlap $q_*$ is chosen
before $D_1$ and before the sequence $(L_n)$, and $D_3$ and $D_4$ do not
depend on the sequence either. Since $L_n\to\infty$, the lemma implies that
$\PP(\mathcal C_n(D_1))\to1$. The event $\mathcal C_n(D_1)$ and the threshold beyond
which the estimates below hold depend on $(L_n)$; the lemma does not provide
one event valid for all sequences at once. For $L_n=n^{1/10000}$,
\cref{b:lem:coverage-extension} contains \cref{inp:ob-coverage}.

\begin{proof}
The proof has three steps. In the first, we show that outside an event of
vanishing probability, every set of Gibbs mass at least $e^{-D_1L_n}$
contains a fraction at least $e^{-D_2L_n}$ of pairs with overlap at least
$q_*$ in absolute value. The idea, from \cite[Section~6]{oai-barriers}, is
that if a set has few such pairs, then a second moment computation shows
that a small rotation of the disorder increases the partition function
restricted to this set by nearly its full Gaussian factor. In the second
step, we show that this gain outweighs the loss caused by shrinking the
original disorder, which the upper energy tail controls, and that the
rotation estimate makes the resulting increase of $\log Z_n^J$ unlikely. In the
third step, we apply the pair bound to the set of exceptional targets
itself.

\paragraph{Constants.}
By \cref{inp:ob-energy-tail}, there exist $e_1\in(0,\beta/2)$ and $c_E>0$
such that the event
\[
 \mathcal E_n=\bigl\{\pi(\{x:H_n^J(x)>ne_1\})\le e^{-c_En}\bigr\}
\]
satisfies $\PP(\mathcal E_n^c)\le e^{-c_En}$ for all sufficiently large $n$.
We define
\[
 g_0=\frac{\beta^2}4-\frac{\beta e_1}2,
\]
which is positive because $e_1<\beta/2$. We fix $q_*\in(0,1)$ with
$\beta^2q_*<g_0/8$. These choices depend only on $\beta$. Given $D_1>0$, we
fix constants $\Lambda$, $D_2$, and $D_3$ such that
\begin{equation}\label{eq:ref-cov-constants}
 \Lambda>\frac{8(D_1+1)}{g_0},\qquad
 D_2>\frac{\beta^2\Lambda}2+1,\qquad
 D_3>D_1+D_2 .
\end{equation}
They depend only on $\beta$ and $D_1$. Now let $(L_n)$ be as in the
statement, write $L=L_n$, and let
\[
 \omega=\omega_n=\frac{\Lambda L_n}n .
\]
Then $\omega_n\to0$ and $n\omega_n=\Lambda L_n\to\infty$. We have
$0<\omega<1/2$ for all sufficiently large $n$; we assume this from now
on.

\paragraph{Step 1: pairs with large overlap in every large set.}
Let $\mathcal V_n$ be the event that some $S\subseteq\Sigma_n$ satisfies
\begin{equation}\label{eq:ref-pair-violation}
 \pi(S)\ge e^{-D_1L}\qquad\text{and}\qquad
 \pi(\cdot\mid S)^{\otimes2}(|R|\ge q_*)<e^{-D_2L}.
\end{equation}
This event is measurable, since it is a finite union of events, each
defined by finitely many continuous functions of the disorder. Off
$\mathcal V_n$, we have
\begin{equation}\label{b:eq:general-coverage-pairs}
 \pi(\cdot\mid S)^{\otimes2}(|R|\ge q_*)\ge e^{-D_2L}
 \qquad\text{for every }S\subseteq\Sigma_n\text{ with }
 \pi(S)\ge e^{-D_1L}.
\end{equation}
We will show that $\PP(\mathcal V_n)\to0$.

Fix $J\in\mathcal E_n\cap\mathcal V_n$ and a set $S$ satisfying
\eqref{eq:ref-pair-violation}. We first remove the configurations of high
energy from $S$. Let
\[
 S'=S\cap\{x:H_n^J(x)\le ne_1\}.
\]
For large $n$, we have $e^{-c_En}\le e^{-D_1L}/2\le\pi(S)/2$, because
$D_1L=o(n)$. This implies that
\[
 \pi(S')\ge\pi(S)-\pi(\{x:H_n^J(x)>ne_1\})\ge\pi(S)/2 .
\]
Restricting the pairs to $S'\times S'$ and dividing by $\pi(S')^2$ instead of
$\pi(S)^2$, we find
\begin{equation}\label{eq:ref-trimmed}
 \pi(\cdot\mid S')^{\otimes2}(|R|\ge q_*)
 \le\Bigl(\frac{\pi(S)}{\pi(S')}\Bigr)^2
   \pi(\cdot\mid S)^{\otimes2}(|R|\ge q_*)
 <4e^{-D_2L}.
\end{equation}

Let $\widetilde J=(\widetilde J_{ij})_{1\le i<j\le n}$ be a standard Gaussian vector independent of
$J$, and define
\[
 W=\sum_{x\in S'}\pi(x\mid S')\,e^{\beta\sqrt\omega H_n^{\widetilde J}(x)},\qquad
 J'=\sqrt{1-\omega}\,J+\sqrt\omega\,\widetilde J .
\]
We write $\PP_{\widetilde J}$ and $\E_{\widetilde J}$ for probability and expectation over $\widetilde J$ with
$J$ fixed. Recall from \eqref{eq:common-covariance} that
$\E H_n^{\widetilde J}(x)H_n^{\widetilde J}(y)=(nR(x,y)^2-1)/2$. Using $\E e^X=e^{\E X^2/2}$ for a
centered Gaussian variable $X$, we find
\[
 \E_{\widetilde J}W=e^{\beta^2\omega(n-1)/4},\qquad
 \frac{\E_{\widetilde J}W^2}{(\E_{\widetilde J}W)^2}
 =\sum_{x,y\in S'}\pi(x\mid S')\,\pi(y\mid S')\exp\Bigl\{\frac{\beta^2\omega}2
     \bigl(nR(x,y)^2-1\bigr)\Bigr\}.
\]
We split the last sum according to whether $|R(x,y)|<q_*$. On the
first part, the exponent is at most
\[
 \Theta=\frac{\beta^2n\omega q_*^2}2 .
\]
On the second part, the exponent is at most $\beta^2n\omega/2=\beta^2\Lambda L/2$,
and we use \eqref{eq:ref-trimmed}. Since $D_2>\beta^2\Lambda/2+1$, this
gives
\[
 \frac{\E_{\widetilde J}W^2}{(\E_{\widetilde J}W)^2}
 \le e^{\Theta}+4e^{(\beta^2\Lambda/2-D_2)L}
 \le e^{\Theta}+4e^{-L}
 \le 2e^{\Theta}
\]
once $L\ge\log4$. The Paley--Zygmund inequality
$\PP(W\ge\E W/2)\ge(\E W)^2/(4\E W^2)$, valid for every nonnegative random
variable $W$ with finite second moment, then implies that
\begin{equation}\label{eq:ref-PZ}
 \PP_{\widetilde J}\bigl(W\ge\E_{\widetilde J}W/2\bigr)\ge\frac{e^{-\Theta}}8 .
\end{equation}

Next, we raise this probability to $1/2$ at the cost of a lower threshold.
Let $f(\widetilde J)=\log W$. Differentiating the finite sum defining $W$, we have
\[
 \frac{\partial f}{\partial\widetilde J_{ij}}=\frac{\beta}{W}\sqrt{\frac\omega n}
   \sum_{x\in S'}\pi(x\mid S')\,e^{\beta\sqrt\omega H_n^{\widetilde J}(x)}x_ix_j .
\]
The right side is $\beta\sqrt{\omega/n}$ times an average of the numbers
$x_ix_j\in\{-1,1\}$, so each coordinate of $\nabla f$ is at most
$\beta\sqrt{\omega/n}$ in absolute value. Summing
over the $n(n-1)/2$ coordinates, we find
$|\nabla f|^2\le\beta^2\omega(n-1)/2=\sigma_W^2$; in particular, $f$ is
$\sigma_W$-Lipschitz. Let $\bar f=\E_{\widetilde J}f$. By the second bound in
\eqref{eq:inp-concentration} (\cref{inp:ob-concentration}), we have $\E_{\widetilde J}\exp\{\theta(f-\bar f)\}\le e^{\theta^2\sigma_W^2/2}$ for every
$\theta\in\R$. The exponential Markov inequality then implies the two
one-sided bounds
\begin{equation}\label{eq:ref-one-sided}
 \PP_{\widetilde J}(f-\bar f\ge a)\le e^{-a^2/(2\sigma_W^2)},\qquad
 \PP_{\widetilde J}(f-\bar f\le-a)\le e^{-a^2/(2\sigma_W^2)}\qquad(a\ge0).
\end{equation}
Note that $W\ge\E_{\widetilde J}W/2$ if and only if $f\ge\log(\E_{\widetilde J}W/2)$. If
$\log(\E_{\widetilde J}W/2)>\bar f$, then \eqref{eq:ref-PZ} and the first bound in
\eqref{eq:ref-one-sided} with $a=\log(\E_{\widetilde J}W/2)-\bar f$ give
\[
 \frac{e^{-\Theta}}8\le\PP_{\widetilde J}\bigl(f\ge\log(\E_{\widetilde J}W/2)\bigr)
 \le\exp\Bigl\{-\frac{(\log(\E_{\widetilde J}W/2)-\bar f)^2}{2\sigma_W^2}\Bigr\},
\]
so that $\bar f\ge\log(\E_{\widetilde J}W/2)-\sigma_W\sqrt{2(\Theta+\log8)}$. This lower
bound for $\bar f$ also holds if $\log(\E_{\widetilde J}W/2)\le\bar f$. The second
bound in \eqref{eq:ref-one-sided} with $a=\sigma_W\sqrt{2\log2}$ then implies
that
\[
 \PP_{\widetilde J}\Bigl(f>\log(\E_{\widetilde J}W/2)-\sigma_W\sqrt{2(\Theta+\log8)}
     -\sigma_W\sqrt{2\log2}\Bigr)\ge\frac12 .
\]
We estimate the loss in the threshold using $\sigma_W^2\le\beta^2n\omega/2$,
$\sqrt{u+v}\le\sqrt u+\sqrt v$, and
$\sqrt{\log8}+\sqrt{\log2}=(1+\sqrt3)\sqrt{\log2}<3$:
\[
 \sigma_W\sqrt{2\Theta}
   +\sigma_W\bigl(\sqrt{2\log8}+\sqrt{2\log2}\bigr)
 \le\frac{\beta^2n\omega q_*}{\sqrt2}+3\beta\sqrt{n\omega} .
\]
Recalling that $\log(\E_{\widetilde J}W/2)=\beta^2\omega(n-1)/4-\log2$, we conclude
that, with conditional probability at least $1/2$,
\begin{equation}\label{b:eq:general-coverage-W}
 \log W\ge\frac{\beta^2\omega(n-1)}4-\beta^2n\omega q_*
   -3\beta\sqrt{n\omega}-\log2 .
\end{equation}
The constants in \eqref{b:eq:general-coverage-W}, and the threshold for $n$
beyond which it holds, do not depend on $J$ or $S$.

\paragraph{Step 2: comparison with the rotation estimate.}
By the linearity of $H_n^J$ in $J$, we have
$H_n^{J'}=\sqrt{1-\omega}\,H_n^J+\sqrt\omega\,H_n^{\widetilde J}$. Because
$1-\sqrt{1-\omega}>0$ and $H_n^J\le ne_1$ on $S'$, we have
\[
 \frac{Z_n^{J'}}{Z_n^J}
 =\pi\Bigl(\exp\bigl\{-\beta\bigl(1-\sqrt{1-\omega}\bigr)H_n^J
    +\beta\sqrt\omega H_n^{\widetilde J}\bigr\}\Bigr)
 \ge\pi(S')\,e^{-\beta(1-\sqrt{1-\omega})ne_1}\,W .
\]
For $0\le\omega\le1$, we have
$1-\sqrt{1-\omega}=\omega/(1+\sqrt{1-\omega})\le\omega/(2-\omega)
\le\omega/2+\omega^2/2$; the last inequality is equivalent to
$\omega(1-\omega)\ge0$. We take logarithms and insert
$\pi(S')\ge e^{-D_1L}/2$, \eqref{b:eq:general-coverage-W}, and
$\beta^2\omega(n-1)/4\ge\beta^2n\omega/4-\beta^2/8$. Recalling that
$n\omega=\Lambda L$, we obtain, with conditional probability at least $1/2$,
\[
 \log Z_n^{J'}-\log Z_n^J
 \ge n\omega\Bigl(\frac{\beta^2}4-\frac{\beta e_1}2-\beta^2q_*\Bigr)
   -D_1L-\mathrm{Err}_n
 =\bigl[(g_0-\beta^2q_*)\Lambda-D_1\bigr]L-\mathrm{Err}_n,
\]
where
\[
 \mathrm{Err}_n=\frac{\beta e_1\Lambda^2L^2}{2n}+3\beta\sqrt{\Lambda L}
   +\frac{\beta^2}8+2\log2 .
\]
By the choices of $q_*$ and $\Lambda$, we have $\beta^2q_*<g_0/8$ and
$D_1<g_0\Lambda/8-1$, so that
\[
 (g_0-\beta^2q_*)\Lambda-D_1>\frac78g_0\Lambda-\frac18g_0\Lambda
 =\frac34g_0\Lambda .
\]
Since $L\to\infty$ and $L/n\to0$, we have $\mathrm{Err}_n=o(L)$; this is
where the two conditions on the scale are used together. We have
$\mathrm{Err}_n\le g_0\Lambda L/4$ for large $n$, and then
\begin{equation}\label{eq:ref-gain}
 \PP_{\widetilde J}\Bigl(\log Z_n^{J'}-\log Z_n^J\ge\frac{g_0\Lambda L}2\Bigr)\ge\frac12
 \qquad\text{for every }J\in\mathcal E_n\cap\mathcal V_n .
\end{equation}
The threshold for $n$ depends only on $\beta$, $D_1$, and the sequence
$(L_n)$.

Since the event in \eqref{eq:ref-gain} does not involve $S$, we need no
measurable choice of $S$ as a function of $J$ (in
\cite[Section~6]{oai-barriers}, the first violating set in a fixed order of
the subsets is chosen; either way works). By Fubini's theorem, the left side
of \eqref{eq:ref-gain} is a measurable function of $J$. Integrating it over
$J$, we obtain
\[
 \frac12\PP(\mathcal E_n\cap\mathcal V_n)
 \le\PP_{J,\widetilde J}\Bigl(\log Z_n^{J'}-\log Z_n^J\ge\frac{g_0\Lambda L}2\Bigr).
\]
The right side is a probability under the joint law of $(J,\widetilde J)$, which is
the setting of \cref{inp:ob-rotation}. Since $0<\omega<1/2$, the last bound
of \cref{inp:ob-rotation}, with $t=\omega$ and $u=g_0\Lambda L/2=g_0n\omega/2$,
bounds it by $\exp\{-u^2/(2\beta^2n\omega)\}=\exp\{-g_0^2\Lambda L/(8\beta^2)\}$.
We conclude that
\[
 \PP(\mathcal V_n)\le\PP(\mathcal E_n^c)+\PP(\mathcal E_n\cap\mathcal V_n)
 \le e^{-c_En}+2\exp\Bigl\{-\frac{g_0^2\Lambda L}{8\beta^2}\Bigr\}
 \longrightarrow0 .
\]
This bound holds simultaneously for all subsets, without a union bound over
them.

\paragraph{Step 3: from pairs to targets.}
Let $\mathcal C_n(D_1)=\mathcal V_n^c$, fix $J\in\mathcal C_n(D_1)$ and
$S\subseteq\Sigma_n$, and let $U_S$ be the set of $x\in S$ with
$\pi(\{z\in S:|R(x,z)|\ge q_*\})<e^{-D_3L}$. Suppose that
$\pi(U_S)\ge e^{-D_1L}$. Since $U_S\subseteq S$, every $x\in U_S$
satisfies $\pi(\{y\in U_S:|R(x,y)|\ge q_*\})<e^{-D_3L}$. This implies that
\begin{align*}
 \pi(\cdot\mid U_S)^{\otimes2}(|R|\ge q_*)
 &=\frac1{\pi(U_S)^2}\sum_{x\in U_S}\pi(x)\,
   \pi(\{y\in U_S:|R(x,y)|\ge q_*\})\\
 &<\frac{e^{-D_3L}}{\pi(U_S)}
 \le e^{-(D_3-D_1)L}<e^{-D_2L},
\end{align*}
where in the last step we used $D_3>D_1+D_2$. Then $U_S$ satisfies
\eqref{eq:ref-pair-violation}, which contradicts $J\notin\mathcal V_n$.
We conclude that $\pi(U_S)<e^{-D_1L}$ for every $S$, which is
\eqref{b:eq:general-coverage}. Finally, for all sufficiently large $n$, the
last display of Step~2 and $L\le n$ give
\[
 \PP\bigl(\mathcal C_n(D_1)^c\bigr)=\PP(\mathcal V_n)
 \le e^{-c_En}+2\exp\Bigl\{-\frac{g_0^2\Lambda L}{8\beta^2}\Bigr\}
 \le3e^{-D_4L},\qquad
 D_4=\min\Bigl\{c_E,\frac{g_0^2\Lambda}{8\beta^2}\Bigr\}.
\]
The constant $D_4$ depends only on $\beta$ and $D_1$. This completes the
proof.
\end{proof}

For $D>0$, we write $c(D)$ for the constant $D_3$ given by
\cref{b:lem:coverage-extension} with $D_1=D$. It depends only on $\beta$ and
$D$.

\section{Parameters and exceptional triples}\label{sec:ref-parameters}

\subsection{Exponents}

We fix exponents $\lambda$ and $\xi$ of the scales such that
\begin{equation}\label{eq:ref-exponents}
 0<\lambda<\xi<\frac14,\qquad 2\xi<3\lambda,
\end{equation}
and we assume the following upper bound on the mean bias
$\Delta_{n,\beta}=P_\beta-p_n(\beta)$: there exists $C_b<\infty$ such that
\begin{equation}\label{eq:ref-mean-input}
 0\le P_\beta-p_n(\beta)\le C_bn^{-3\lambda}
 \qquad\text{for all sufficiently large }n.
\end{equation}
For \cref{thm:barrier-730}, we take $\lambda=7/30$ and $\xi=5/21$, and
\eqref{eq:ref-mean-input} is \cref{inp:hb-rate}, with exponent $7/10=3\lambda$.
Then \eqref{eq:ref-exponents} holds because $7/30<5/21<1/4$ and
$2\xi=10/21<7/10=3\lambda$. The barrier exponent is one
third of the exponent of the mean-bias input. \Cref{tab:ref-exponents}
lists the powers of $n$ that appear in this section and in
\cref{sec:induction,sec:barrier-proof}; the error terms of the static
estimates are tabulated in \cref{sec:static}.

\subsection{The fixed parameters}

Let $q_*=q_*(\beta)$ be the overlap of \cref{b:lem:coverage-extension}.
Recall from \cref{lem:st-fixed} that for each fixed large-overlap cutoff
$b_*\in(0,1)$, there exists a small-overlap cutoff $\underline b(b_*)>0$, depending only
on $\beta$ and $b_*$, such that the three-replica pressure deficit
\eqref{b:eq:fixed-pressure-gap} holds for all feasible overlap matrices
$D(r,q,\delta)$ with $q\ge b_*$, $|r+\delta|\le\underline b(b_*)$,
$|r-\delta|\le\underline b(b_*)$, and $|q-\Gamma(q)|\le\nu_{\rm fix}(b_*)$. We fix signed-overlap cutoffs
$b_1$ and $b_2$ with
\[
 0<b_1<1,\quad b_1\le\underline b(q_*),\qquad
 0<b_2<b_1,\quad b_2\le\underline b(b_1),
\]
and then $t_0>0$ with
\begin{equation}\label{b:dyn:fixed-cutoffs}
 4t_0<\min\{q_*,b_2,b_1^2/2\},\qquad 3t_0<q_\beta .
\end{equation}
In \cref{sec:induction}, the thresholds lie in $[t_0,2t_0]$, and the cells
that we define below use the cutoff $3t_0$. We use the first condition in
\eqref{b:dyn:fixed-cutoffs} to keep these values below $q_*$ and $b_2$, with
room for the error terms of \cref{sec:induction}. It also gives
$3t_0<b_1^2/2$, which we use there to bound the number of deleted
references. The second condition ensures that the intervals $J_\beta$ and
$I_\beta$ chosen next lie inside $(0,q_\beta)$, where $\Gamma(t)=t$ by
\eqref{b:eq:identity-interval}. We fix compact intervals $J_\beta$ and
$I_\beta$ such that
\begin{equation}\label{eq:ref-intervals}
 [t_0,2t_0]\subseteq\operatorname{int}J_\beta,\qquad
 J_\beta\subseteq\operatorname{int}I_\beta,\qquad
 I_\beta\subseteq(t_0/2,3t_0).
\end{equation}
The centers of all three-replica constraints met in \cref{sec:induction}
lie in $J_\beta$. The local comparison of
\cref{sec:fractional,sec:static} is applied only at centers in $J_\beta$,
with its local cutoff chosen so small that all its windows lie in
$I_\beta$. Let $m_\beta>0$ be the lower
bound of \eqref{b:eq:main-analytic-constants} for
$\alpha(r)^2\beta^6\E\chi(r,Y_r)^3$ on $I_\beta$; it is positive because
$I_\beta$ is a compact subset of $(0,q_\beta)$.

We fix the band of the local comparison of \cref{sec:cubic} at
$\vartheta=1/2$ and $E=1/100$. Here $\vartheta$ is the ratio $(q-r)/h$ of
that comparison and $E\rho_n\le|\delta|\le2E\rho_n$ is its range of
imbalances. At these values \eqref{b:eq:os-band-checks} holds and
$G(\vartheta,E)>1/3$. We fix the pressure coefficient $g$, a fraction
$\varkappa\in(0,1)$, and the scale constant $R>0$ with
\begin{align}
 0<g&<m_\beta E^2G(\vartheta,E),\label{b:dyn:finite-band-g}\\
 R^3&>\frac{3C_b}{\varkappa g},\label{b:dyn:R-condition}
\end{align}
where $C_b$ is the constant of \eqref{eq:ref-mean-input}. Here $g$ is the
coefficient of $\rho_n^3$ in the local fractional gap
\eqref{b:eq:local-fractional-gap}, and $\varkappa$ is the fraction of this gap that
is kept in the probability bounds of \cref{sec:static}. The scales are
\begin{equation}\label{b:dyn:scales}
 \rho_n=Rn^{-\lambda},\qquad \varepsilon_n=E\rho_n,\qquad
 \eta_n=n^{-\xi},
\end{equation}
which are the scales of \cref{rem:st-scales}. With these choices, the data
of the simultaneous estimate in \cref{sec:st-families} are the intervals
$J_\beta\subseteq\operatorname{int}I_\beta$, the large-overlap cutoffs
$q_*$ and $b_1$, the band $(\vartheta,E)$, the coefficient $g$, and the
fraction $\varkappa$. The remaining constants of the static estimates (the local
cutoff, the tolerance $\nu$ of the pair exclusions, the pressure deficits of
the separated-overlap, fixed-scale, and pair regimes, the coefficient of the
fractional tilt, and the denominator margin \eqref{b:eq:denominator-margin})
are fixed in \cref{sec:fractional,sec:static} with positive slack. They
depend only on $\beta$ and on the parameters above, and not on $n$.

Finally, we fix integers
\begin{equation}\label{b:dyn:BH}
 H=\lceil2/b_1^2\rceil,\qquad B\ge H+2 .
\end{equation}
Here $H$ bounds the number of references deleted at the start of a block in
\cref{sec:induction}, and $B$ is the number of levels in a block; $B$ is
chosen in \cref{sec:barrier-proof}. After $B$, the constants
$c_s,C_s$ of the schedule \eqref{eq:ref-schedule} below are fixed, and then
the coefficient $A>0$ of the logarithmic time scale
\[
 L_n=An^\lambda
\]
is chosen in \cref{sec:barrier-proof}, with $A=A_\beta$ for \cref{thm:barrier-730}; the time horizon is $e^{L_n}$. The sequence $L_n$ tends to infinity
and is $o(n)$, because $0<\lambda<1$; in particular,
\cref{b:lem:coverage-extension} applies to it.

\begin{table}[t]
\centering\small
\caption{The fixed parameters of the dynamical argument, in the order in
which they are chosen. Every entry is fixed before $n\to\infty$; each
depends only on $\beta$ and on the entries above it.}
\label{tab:ref-parameters}
\begin{tabular}{@{}>{$}l<{$}p{7.4cm}p{4.0cm}@{}}
\toprule
\text{Parameter} & Choice & Reference\\
\midrule
q_\beta & right end of $\operatorname{supp}\mu_\beta=[0,q_\beta]$ & \cref{inp:full-rsb}\\
q_* & overlap of the overlap-neighborhood mass bound & \cref{b:lem:coverage-extension}\\
b_1,\ b_2 & $b_1\le\underline b(q_*)$, $b_2\le\underline b(b_1)$, $0<b_2<b_1<1$ & \eqref{b:eq:fixed-pressure-gap}\\
t_0 & $4t_0<\min\{q_*,b_2,b_1^2/2\}$, $3t_0<q_\beta$ & \eqref{b:dyn:fixed-cutoffs}\\
J_\beta,\ I_\beta & compact intervals around $[t_0,2t_0]$ & \eqref{eq:ref-intervals}\\
m_\beta & lower bound for the scalar coefficient on $I_\beta$ & \eqref{b:eq:main-analytic-constants}\\
\vartheta,\ E & $\vartheta=1/2$, $E=1/100$ & \eqref{b:eq:os-band-checks}\\
g & $0<g<m_\beta E^2G(\vartheta,E)$ & \eqref{b:dyn:finite-band-g}\\
\varkappa & $0<\varkappa<1$ & \cref{sec:static}\\
C_b & constant of the mean-bias input; for \cref{thm:barrier-730}, that of \cref{inp:hb-rate} & \eqref{eq:ref-mean-input}\\
R & $R^3>3C_b/(\varkappa g)$ & \eqref{b:dyn:R-condition}\\
\text{static constants} & local cutoff, $\nu$, deficits, fractional tilt, denominator margin & \cref{sec:fractional,sec:static}\\
H & $H=\lceil2/b_1^2\rceil$ & \eqref{b:dyn:BH}\\
\text{threshold constants} & widths and spacing of the threshold intervals & \cref{sec:induction}\\
B & integer, $B\ge H+2$, block length & \cref{sec:barrier-proof}\\
c_s,\ C_s & schedule of list sizes, $0\le s<B$ & \eqref{eq:ref-schedule}\\
A & coefficient of $L_n=An^\lambda$ & \cref{sec:barrier-proof}\\
\bottomrule
\end{tabular}
\end{table}

We note that the quantities that depend on $n$ are the scales $\rho_n$,
$\varepsilon_n$, and $\eta_n$ of \eqref{b:dyn:scales}, the logarithmic time scale $L_n$, the
number $N_{\rm lev}$ of levels (fixed in \cref{sec:induction}), the list
sizes $K_s$ of \eqref{eq:ref-schedule}, and the total number $N$ of list
entries. We never let a fixed parameter vary with $n$, and the threshold
$\tau$ for the total variation distance does not enter until the last step
of \cref{sec:barrier-proof}.

\begin{table}[t]
\centering\small
\caption{Powers of $n$ at the scales \eqref{b:dyn:scales} and the logarithmic
time scale $L_n=An^\lambda$, for general exponents and for \cref{thm:barrier-730}.}
\label{tab:ref-exponents}
\begin{tabular}{@{}>{$}l<{$}>{$}l<{$}>{$}l<{$}p{3.9cm}@{}}
\toprule
\text{Quantity} & \text{General} & \text{\cref{thm:barrier-730}} & Role\\
\midrule
P_\beta-p_n(\beta) & O(n^{-3\lambda}) & O(n^{-7/10}) & mean-bias input\\
\rho_n & Rn^{-\lambda} & Rn^{-7/30} & overlap gap\\
\eta_n & n^{-\xi} & n^{-5/21} & fixed-scale tolerance\\
L_n & An^{\lambda} & A_\beta n^{7/30} & logarithmic time scale\\
n\rho_n^3 & R^3n^{1-3\lambda} & R^3n^{3/10} & narrow triples\\
n\eta_n^2 & n^{1-2\xi} & n^{11/21} & fixed-scale triples\\
n\rho_n^3/L_n & (R^3/A)n^{1-4\lambda} & (R^3/A_\beta)n^{1/15} & \eqref{eq:ref-absorb}\\
n\eta_n^2/L_n & n^{1-2\xi-\lambda}/A & n^{61/210}/A_\beta & \eqref{eq:ref-absorb}\\
\eta_n/\rho_n & n^{\lambda-\xi}/R & n^{-1/210}/R & \cref{sec:induction}\\
\bottomrule
\end{tabular}
\end{table}

\subsection{Exceptional triples}

The overlap constraints that we use in \cref{sec:induction} are of two
kinds. The first kind states that when $x$ and ${\widetilde v}$ have a large overlap,
their small overlaps with a third configuration $v$ are nearly equal. For a
large-overlap cutoff $b_*\in(0,1)$ and a small-overlap cutoff $b'\in(0,1)$, we
define
\begin{equation}\label{eq:ref-fixed-violation}
 \begin{aligned}
 \mathcal B_n(b_*,b')=\bigl\{(v,x,{\widetilde v})\in\Sigma_n^3:{}
   &R(x,{\widetilde v})\ge b_*,\ |R(v,x)|\le b',\ |R(v,{\widetilde v})|\le b',\\
   &|R(v,x)-R(v,{\widetilde v})|>2\eta_n\bigr\}.
 \end{aligned}
\end{equation}
This is the set excluded in \cite[Lemma~4.1]{oai-barriers}, with the
tolerance $2n^{-1/100}$ there replaced by $2\eta_n$. We bound its Gibbs
probability at this tolerance in \cref{sec:static} through the pressure
estimate \eqref{b:eq:fixed-pressure-gap}, and not through
\cite[Lemma~4.1]{oai-barriers}. A triple
$(v,x,{\widetilde v})\notin\mathcal B_n(b_*,b')$ such that $R(x,{\widetilde v})\ge b_*$,
$|R(v,x)|\le b'$, and $|R(v,{\widetilde v})|\le b'$ satisfies
$|R(v,x)-R(v,{\widetilde v})|\le2\eta_n$.

The second kind concerns overlaps near a level $r$ in the compact interval
$J_\beta$. For $(v,x,{\widetilde v})\in\Sigma_n^3$, we write
\begin{equation}\label{b:dyn:narrow-coordinates}
 r=\frac{R(v,x)+R(v,{\widetilde v})}2,\qquad
 \delta=\frac{R(v,x)-R(v,{\widetilde v})}2,\qquad
 q=R(x,{\widetilde v}),
\end{equation}
so that the Gram matrix of $(v,x,{\widetilde v})$ is the matrix $D(r,q,\delta)$ of
\cref{sec:static}, with off-diagonal entries $r+\delta$, $r-\delta$, and $q$.
Let $\mathcal N_n$ be the set of triples with
\[
 r\in J_\beta,\qquad \varepsilon_n\le|\delta|\le2\varepsilon_n,\qquad
 q\ge r+\rho_n .
\]
A triple outside $\mathcal N_n$ with $r\in J_\beta$ and
$\varepsilon_n\le|\delta|\le2\varepsilon_n$ satisfies $R(x,{\widetilde v})<r+\rho_n$: if
the overlaps of $v$ with $x$ and with ${\widetilde v}$ lie near a common level $r$ and
differ by an amount of order $\varepsilon_n$, then the overlap of $x$ and ${\widetilde v}$
cannot exceed $r$ by $\rho_n$. The set $\mathcal N_n$ corresponds to the set
excluded in \cite[Lemma~7.1]{oai-barriers}. There, the gap is
$\rho=n^{-1/10000}$, the imbalance is of order $\rho^{20}$, and the centers
are restricted by the condition $\mu_\beta([r-\rho^8,r+\rho^8])\le\rho^4$ on
the Parisi measure. Here the gap $\rho_n$ and the imbalance
$\varepsilon_n=E\rho_n$ are of the same order, and there is no condition on
the Parisi measure, because the local comparison of \cref{sec:fractional} is
uniform over centers in $J_\beta$. The exceptional set is
\begin{equation}\label{eq:ref-exceptional}
 \mathcal B_n=\mathcal N_n\cup\mathcal B_n(q_*,b_1)\cup\mathcal B_n(b_1,b_2).
\end{equation}
It is a deterministic subset of $\Sigma_n^3$, and it is not closed under
permutations or spin flips of the coordinates. In the dynamical argument, we
evaluate $\mathcal B_n$ on triples in several orders and orientations; since
$\pi^{\otimes3}$ is invariant under these operations, each of them costs only
a fixed factor in \cref{lem:ref-families}.

The Gibbs probability of $\mathcal B_n$ is controlled in \cref{sec:static},
where different parts of it are handled by different estimates. The set
$\mathcal N_n$ is the set of narrow triples of \cref{sec:st-families} for the
interval of centers $J_\beta$. By \cref{lem:st-cover}, a narrow triple either
has its pair $(x,{\widetilde v})$ excluded by the pair deficit
\eqref{b:eq:pair-static-gap}, or its Gram matrix lies in the local family,
for which we use the positive fractional moment
\eqref{b:eq:local-fractional-gap}, or in the separated-overlap family, for
which we use the stronger pressure deficit \eqref{b:eq:remote-pressure-gap}
and ordinary Gaussian concentration. The sets $\mathcal B_n(q_*,b_1)$ and
$\mathcal B_n(b_1,b_2)$ are contained in the set $\mathcal W_n$ of
fixed-scale triples of \cref{sec:st-families}, which is controlled by
\eqref{b:eq:fixed-pressure-gap} and the pair deficit. We need
no fractional moment bound at a distant overlap $q$.

\begin{lemma}[One disorder event for the exceptional triples]
\label{lem:ref-single-event}
Let $(\rho_n)$ and $(\eta_n)$ be deterministic positive sequences with
$\rho_n\to0$ that satisfy \eqref{eq:st-S1}--\eqref{eq:st-S5}, and let
$\varepsilon_n=E\rho_n$. Let
$\mathcal B_n$ be the set \eqref{eq:ref-exceptional} defined with these
sequences, and let $\mathcal A_n$ be the disorder event of
\cref{b:lem:static-event} for these sequences and for the data of
\cref{sec:st-families} fixed in \cref{tab:ref-parameters} (in particular the
intervals $J_\beta$ and $I_\beta$ and the large-overlap cutoffs $q_*$ and
$b_1$). Then $\PP(\mathcal A_n)\to1$, and there exist constants $C,c>0$ such
that, for all sufficiently large $n$ and every $J\in\mathcal A_n$,
\begin{equation}\label{b:dyn:uniform-gibbs}
 \pi^{\otimes3}(\mathcal B_n)\le q_n,\qquad
 q_n=C(n+1)^3\bigl(e^{-cn\rho_n^3}+e^{-cn\eta_n^2}+e^{-cn}\bigr).
\end{equation}
\end{lemma}

\begin{proof}
By \cref{b:lem:static-event}, we have $\PP(\mathcal A_n^c)\to0$. Moreover,
for large $n$ and every $J\in\mathcal A_n$,
\eqref{b:eq:conditional-static-errors} bounds the Gibbs probability of the set
$\mathcal N_n$ of narrow triples of \cref{sec:st-families} by
$C(n+1)^3e^{-cn\rho_n^3}$, and that of the set $\mathcal W_n$ of fixed-scale
triples of \cref{sec:st-families} by $C(n+1)^3e^{-cn\eta_n^2}$.

It remains to compare $\mathcal B_n$ with these two sets, where our triples
$(v,x,{\widetilde v})$ are the triples $(v,x,y)$ of \cref{sec:static}. The set
$\mathcal N_n$ in \eqref{eq:ref-exceptional} is the set of narrow triples for
the interval of centers $J_\beta$. Let $(v,x,{\widetilde v})\in\mathcal B_n(b_*,b')$ with
$(b_*,b')=(q_*,b_1)$ or $(b_*,b')=(b_1,b_2)$. Then $R(x,{\widetilde v})\ge b_*$,
$|R(v,x)|\le b'\le\underline b(b_*)$, $|R(v,{\widetilde v})|\le b'\le\underline b(b_*)$,
and $|R(v,x)-R(v,{\widetilde v})|>2\eta_n$. Since $b_*$ is one of the cutoffs $q_*$ and
$b_1$, the triple satisfies the conditions that define $\mathcal W_n$. This
shows that $\mathcal B_n(q_*,b_1)\cup\mathcal B_n(b_1,b_2)\subseteq\mathcal W_n$.
Summing the two bounds of the previous paragraph, we obtain
$\pi^{\otimes3}(\mathcal B_n)\le C(n+1)^3(e^{-cn\rho_n^3}+e^{-cn\eta_n^2})$,
which implies \eqref{b:dyn:uniform-gibbs}.
\end{proof}

We apply \cref{lem:ref-single-event} with the scales \eqref{b:dyn:scales}.
Since $\lambda>0$, we have $\rho_n=Rn^{-\lambda}\to0$. These scales also satisfy
\eqref{eq:st-S1}--\eqref{eq:st-S5}, as in
\cref{rem:st-scales}(b). Indeed, by \eqref{eq:ref-mean-input} and
\eqref{b:dyn:R-condition}, we have
$\Delta_{n,\beta}/\rho_n^3\le C_b/R^3<\varkappa g/3$ for large $n$, which is
\eqref{eq:st-S1}. Further, $n^{7/15}/(n^2\rho_n^6)=R^{-6}n^{6\lambda-23/15}$
tends to zero because $6\lambda<3/2<23/15$; the quantities
$n\rho_n^4=R^4n^{1-4\lambda}$ and $n\eta_n^4=n^{1-4\xi}$ grow faster than
$\log n$ because $\lambda,\xi<1/4$; and $\rho_n^3/\eta_n^2=R^3n^{2\xi-3\lambda}$
tends to zero because $2\xi<3\lambda$. From now on, $\mathcal B_n$,
$\mathcal A_n$, and $q_n$ denote the objects of \cref{lem:ref-single-event} at
the scales \eqref{b:dyn:scales}. In \cref{sec:bp-strong}, the lemma is
applied with $\eta_n=\rho_n/1000$ instead.

\begin{remark}[Two kinds of error]\label{rem:ref-two-errors}
In the proof of \cref{thm:barrier-730}, we use two kinds of
small quantities, and we never mix them. The bound
\eqref{b:eq:static-event-failure} on the disorder probability
$\PP(\mathcal A_n^c)$ includes a polynomial term from Chebyshev's inequality,
and therefore it does not remain small after multiplication by a factor
$e^{CL_n}$. For $\lambda=7/30$ and $\xi=5/21$, the terms
$e^{-cn\rho_n^4}$ and $e^{-cn\eta_n^4}$ in it, which are
$e^{-cR^4n^{1/15}}$ and $e^{-cn^{1/21}}$, also decay more slowly than
$e^{-CL_n}=e^{-CAn^{7/30}}$ for every fixed $C>0$. We remove this disorder probability,
together with the probabilities of the complements of the events of
\cref{b:lem:coverage-extension}, once, before counting reference
configurations or path positions. The conditional Gibbs bound $q_n$ in
\eqref{b:dyn:uniform-gibbs}, in contrast, holds for every
$J\in\mathcal A_n$, and it absorbs every factor
$e^{CL_n}$. \Cref{tab:ref-exponents} shows that the exponents
$n\rho_n^3$, $n\eta_n^2$, and $n$ exceed $L_n$ by a power of $n$. We verify
in the proof of \cref{lem:ref-families} that
\begin{equation}\label{eq:ref-absorb}
 \lim_{n\to\infty}e^{CL_n}q_n=0\qquad\text{for every fixed }C>0 .
\end{equation}
In \cref{lem:ref-families}, only $q_n$ is multiplied by the number of
reference configurations. The deterministic set $\mathcal B_n$ contains every
exceptional triple that can arise for the allowed thresholds and paths, and \eqref{b:dyn:uniform-gibbs} bounds its mass under
$\pi^{\otimes3}$ by $q_n$. We impose no condition
on the disorder that depends on the thresholds of \cref{sec:induction}.
\end{remark}

\section{Fixed families of reference configurations}\label{sec:ref-families}

The induction of \cref{sec:induction} selects references from lists of
independent Gibbs samples, organized in blocks of $B$ consecutive levels.
This is the construction of \cite[Section~8.3]{oai-barriers}. The selection
cell at a level is determined by the references chosen at the preceding
levels of the same block. The list sizes are chosen recursively over the $B$
positions of a block, so that the depth of the recursion is the fixed number
$B$ even though the number of levels tends to infinity. In
\cite[Section~8.3]{oai-barriers}, the required properties of the lists are
proved with probability tending to one jointly over the disorder, the lists,
and the path. The proof there uses union bounds over path positions and over
triples, and the bounds for the triples are averaged over the disorder. At
our scales, the
averaged bound available from \cref{b:lem:static-event} contains the disorder
failure probability of \cref{rem:ref-two-errors}, and this term does not
remain small after multiplication by the number $e^{O(L_n)}$ of reference
configurations. We therefore fix the
disorder first, in the event of \cref{lem:ref-single-event}, and use only
the conditional bound \eqref{b:dyn:uniform-gibbs}. At such a disorder, we fix
one realization of the lists, together with a set $V$ of exceptional
targets, before any path is considered.

\subsection{Lists, cells, and list sizes}

Let $N_{\rm lev}=N_{\rm lev}(n)\ge1$ be the number of levels; it is fixed in
\cref{sec:induction}. Level $j\in\{0,\ldots,N_{\rm lev}-1\}$ belongs to block
$\lfloor j/B\rfloor$ and has position $s(j)=j-B\lfloor j/B\rfloor$ in its
block. The list at level $j$ is an ordered tuple
\[
 \mathcal F_j=(y_{j,1},\ldots,y_{j,K_{s(j)}})\in\Sigma_n^{K_{s(j)}},
\]
whose coordinates we call the entries of $\mathcal F_j$; the sizes $K_s$ are
defined in \eqref{eq:ref-schedule} below. We write
\[
 \mathcal F_j^\pm=\bigl\{\sigma y:\sigma\in\{-1,1\},\ y\text{ an entry of }
   \mathcal F_j\bigr\}
\]
for the signed list. The family of reference configurations is
$\mathcal F=(\mathcal F_0,\ldots,\mathcal F_{N_{\rm lev}-1})$, and
$N=\sum_{j<N_{\rm lev}}K_{s(j)}$ is its total number of entries.

For a level $j$, let $\mathcal T_j$ be the set of tuples $T=(u_i)_{i\in\mathcal I}$,
where $\mathcal I$ is a set of levels of the block of $j$ that precede $j$, and $u_i$
is an entry of $\mathcal F_i$ for each $i\in\mathcal I$. The cell of $T$ is
\begin{equation}\label{eq:ref-cell}
 S_T=\bigl\{x\in\Sigma_n:|R(u,x)|\le3t_0\text{ for every entry }u
   \text{ of }T\bigr\}.
\end{equation}
The empty tuple has cell $S_\emptyset=\Sigma_n$, and it is the only element
of $\mathcal T_j$ when $s(j)=0$. Spin flips of the entries do not change the
cells. The references selected in \cref{sec:induction} at the levels of a
block that precede $j$ form a tuple in $\mathcal T_j$, and the cells used
there are among the cells $S_T$.

The list sizes are chosen by the following schedule over positions
$s=0,\ldots,B-1$. Recall that $c(D)$ denotes the constant $D_3$ of
\cref{b:lem:coverage-extension} for $D_1=D$. We set $C_0=0$ and, recursively,
\begin{equation}\label{eq:ref-schedule}
 c_s=c(C_s+5),\qquad C_{s+1}=C_s+c_s+2,\qquad
 K_s=\bigl\lceil e^{(c_s+1)L_n}\bigr\rceil .
\end{equation}
We have $C_s=\sum_{s'<s}(c_{s'}+2)$. The constants $c_s$ and $C_s$ depend only
on $\beta$ and $s$, and only $K_s$ depends on $n$ and $A$. For $D>0$, let
$\mathcal C_n(D)$ be the event of \cref{b:lem:coverage-extension} for the
sequence $(L_n)$. On $\mathcal C_n(D)$, \eqref{b:eq:general-coverage} reads
\begin{equation}\label{b:dyn:coverage}
 \pi\bigl(\bigl\{x\in S:\pi(\{y\in S:|R(x,y)|\ge q_*\})<e^{-c(D)L_n}\bigr\}
   \bigr)<e^{-DL_n}\qquad\text{for every }S\subseteq\Sigma_n .
\end{equation}
Only $B$ of these events are needed. We define the disorder event
\begin{equation}\label{eq:ref-good-event}
 \mathcal G_n=\mathcal A_n\cap\bigcap_{s=0}^{B-1}\mathcal C_n(C_s+5).
\end{equation}
It is measurable, and $\PP(\mathcal G_n)\to1$ by
\cref{lem:ref-single-event,b:lem:coverage-extension}, since it is the
intersection of a fixed number of events whose probabilities tend to one.

\subsection{The family}

\begin{lemma}[Fixed families of reference configurations]
\label{lem:ref-families}
Let $B\ge1$ be a fixed integer, let $(L_n)$ be a deterministic sequence with
$L_n\to\infty$ and $L_n/n\to0$, and let $K_s$ and $\mathcal G_n$ be given by
\eqref{eq:ref-schedule} and \eqref{eq:ref-good-event}. Let
$N_{\rm lev}=N_{\rm lev}(n)\ge1$ be integers with
$N_{\rm lev}\le e^{L_n}$, and assume that $q_n$ satisfies
\eqref{eq:ref-absorb}. Then, for all sufficiently large $n$ and every
$J\in\mathcal G_n$, there exist lists
$\mathcal F_j\in\Sigma_n^{K_{s(j)}}$, $0\le j<N_{\rm lev}$, and a set
$V\subseteq\Sigma_n$ with
\begin{equation}\label{b:dyn:fixed-bank}
 \pi(V)\le3e^{-4L_n},
\end{equation}
such that the following four properties hold.
\begin{enumerate}[label=\textup{(\roman*)}]
\item \emph{Overlap neighborhoods in every cell of the current block.} For
every level $j$,
every $T\in\mathcal T_j$, and every $x\in S_T\setminus V$, there exists an
entry $y$ of $\mathcal F_j$ with $y\in S_T$ and $|R(x,y)|\ge q_*$.
\item \emph{Comparison across the overlap $q_*$.} For every
$x\in\Sigma_n\setminus V$, all levels $i\ne j$, and all
$v\in\mathcal F_i^\pm$ and ${\widetilde v}\in\mathcal F_j^\pm$ with
\[
 R(x,{\widetilde v})\ge q_*,\qquad |R(v,x)|\le b_1,\qquad |R(v,{\widetilde v})|\le b_1,
\]
we have $|R(v,x)-R(v,{\widetilde v})|\le2\eta_n$.
\item \emph{Comparison across the overlap $b_1$ for triples of entries.}
For all pairwise distinct levels $i$, $a$, and $j$, and all
$v\in\mathcal F_a^\pm$, $y\in\mathcal F_i^\pm$, and ${\widetilde v}\in\mathcal F_j^\pm$
with
\[
 |R(y,{\widetilde v})|\ge b_1,\qquad |R(v,y)|\le b_2,\qquad |R(v,{\widetilde v})|\le b_2,
\]
we have $|R(v,y)-\sigma R(v,{\widetilde v})|\le2\eta_n$, where $\sigma\in\{-1,1\}$ is the
sign of $R(y,{\widetilde v})$. In particular,
$\bigl||R(v,y)|-|R(v,{\widetilde v})|\bigr|\le2\eta_n$.
\item \emph{No exceptional narrow triple.} For every
$x\in\Sigma_n\setminus V$, all levels $i\ne j$, and all
$v\in\mathcal F_i^\pm$ and ${\widetilde v}\in\mathcal F_j^\pm$, we have
$(v,x,{\widetilde v})\notin\mathcal N_n$. That is, if the coordinates
\eqref{b:dyn:narrow-coordinates} of $(v,x,{\widetilde v})$ satisfy $r\in J_\beta$ and
$\varepsilon_n\le|\delta|\le2\varepsilon_n$, then $R(x,{\widetilde v})<r+\rho_n$.
\end{enumerate}
If $L_n=An^\lambda$ for a fixed $A>0$, with $\lambda$ and $\xi$ as in
\eqref{eq:ref-exponents}, then $L_n\to\infty$, $L_n/n\to0$, and
\eqref{eq:ref-absorb} holds.
\end{lemma}

The threshold for $n$ depends on $\beta$, on the parameters of
\cref{tab:ref-parameters}, and on the sequences $(L_n)$ and
$(N_{\rm lev})$, but not on $J$. In \cref{sec:induction,sec:barrier-proof},
the lemma is applied with $L_n=An^\lambda$ and the number of levels fixed
there, which is at most a constant times $n^\lambda$.

\begin{proof}
The proof is an expectation argument over independent samples of the lists.
We bound the expected Gibbs mass of the targets missed by the lists, that is,
of the configurations $x$ such that, at some level $j$, some cell $S_T$
containing $x$ contains no entry $y$ of $\mathcal F_j$ with
$|R(x,y)|\ge q_*$. We also bound the probability that three list entries
form an exceptional triple, and the expected Gibbs mass of the configurations
that form an exceptional triple with two list entries. We then choose a realization at which the sum of the Gibbs mass of the union of these two sets and the indicator of that event is at most its mean.

\paragraph{Step 1: random lists.}
Fix $n$ and $J\in\mathcal G_n$, and write $L=L_n$. We let the entries
$y_{j,k}$, $0\le j<N_{\rm lev}$, $1\le k\le K_{s(j)}$, be independent random
configurations with law $\pi$, defined on an auxiliary probability space with
probability $\PP_{\mathcal F}$ and expectation $\E_{\mathcal F}$. The
disorder is fixed throughout, and the space $\Sigma_n^N$ of realizations is
finite.

\paragraph{Step 2: targets missed by the lists.}
For a level $j$ and $T\in\mathcal T_j$, let $U_{j,T}$ be the set of
$x\in S_T$ such that no entry $y$ of $\mathcal F_j$ satisfies $y\in S_T$ and
$|R(x,y)|\ge q_*$, and let $V_1=\bigcup_j\bigcup_{T\in\mathcal T_j}
U_{j,T}$. Fix $j$, let $s=s(j)$, and condition on the lists at the levels of
the block of $j$ that precede $j$. This fixes $\mathcal T_j$ and the cells
$S_T$, and leaves the $K_s$ entries of $\mathcal F_j$ independent with law
$\pi$. For $T\in\mathcal T_j$ and $x\in S_T$, let
$a_T(x)=\pi(\{y\in S_T:|R(x,y)|\ge q_*\})$. Then
\[
 \E_{\mathcal F}\bigl[\pi(U_{j,T})\bigm|\text{earlier lists}\bigr]
 =\sum_{x\in S_T}\pi(x)\bigl(1-a_T(x)\bigr)^{K_s}.
\]
We split the sum according to whether $a_T(x)<e^{-c_sL}$. Recall that
$J\in\mathcal C_n(C_s+5)$ and $c_s=c(C_s+5)$. By \eqref{b:dyn:coverage} with
$S=S_T$, the first group has Gibbs mass less than $e^{-(C_s+5)L}$. For the
second group, $(1-a_T(x))^{K_s}\le\exp\{-a_T(x)K_s\}\le
\exp\{-e^{-c_sL}e^{(c_s+1)L}\}=e^{-e^L}$. The number of tuples is
\[
 |\mathcal T_j|\le\prod_{s'<s}\bigl(K_{s'}+1\bigr)
 \le\prod_{s'<s}e^{(c_{s'}+2)L}=e^{C_sL},
\]
where we used $K_{s'}+1\le e^{(c_{s'}+1)L}+2\le e^{(c_{s'}+2)L}$, valid
once $e^L\ge3$. Summing over $T\in\mathcal T_j$ and taking expectations, we
obtain
\[
 \E_{\mathcal F}\,\pi\Bigl(\bigcup_{T\in\mathcal T_j}U_{j,T}\Bigr)
 \le e^{C_sL}\bigl(e^{-(C_s+5)L}+e^{-e^L}\bigr)\le2e^{-5L},
\]
where the last inequality holds once $e^L\ge(C_s+5)L$ for every $s<B$.
Summing over the $N_{\rm lev}\le e^L$ levels, we find
\begin{equation}\label{eq:ref-miss-mass}
 \E_{\mathcal F}\,\pi(V_1)\le2e^{-4L}.
\end{equation}

\paragraph{Step 3: exceptional triples.}
Let $c_{\max}=\max_{s<B}c_s$. Using $N_{\rm lev}\le e^L$ and
$K_s\le e^{(c_s+2)L}$, we find $N\le e^{C_NL}$ with $C_N=c_{\max}+3$.

Let $\mathcal O$ be the event that there exist entries $y_1$, $y_2$, and
$y_3$ of lists at three distinct levels, and signs
$\sigma_1,\sigma_2,\sigma_3\in\{-1,1\}$, such that
$(\sigma_1y_1,\sigma_2y_2,\sigma_3y_3)\in\mathcal B_n$. Fix a choice of three entries, specified by their levels and
indices, and of signs. The three entries are independent with law $\pi$,
because they belong to distinct levels, and $\pi$ is invariant under spin
flips. The triple therefore has law $\pi^{\otimes3}$, and it lies in
$\mathcal B_n$ with probability at most $q_n$ by
\eqref{b:dyn:uniform-gibbs}. There are at most $N^3$ ordered choices of
three entries and $8$ choices of signs, which gives
\[
 \PP_{\mathcal F}(\mathcal O)\le8N^3q_n .
\]
The configurations themselves need not be distinct; only their sampling
indices are.

Let $V_2$ be the set of $x\in\Sigma_n$ for which there exist entries
$y$ and $y'$ of lists at two distinct levels, signs $\sigma,\sigma',
\sigma''\in\{-1,1\}$, and an ordering of $\sigma y$, $\sigma'y'$, and
$\sigma''x$ as a triple, such that this triple lies in $\mathcal B_n$. Let
$x$ be a configuration with law $\pi$, independent of the lists. For a fixed
ordered pair of entries at distinct levels, fixed signs, and a fixed
ordering, the resulting triple has law $\pi^{\otimes3}$, and it lies in
$\mathcal B_n$ with probability at most $q_n$. There are at most $N^2$
ordered pairs of entries, $8$ choices of signs, and $6$ orderings, and a
union bound yields
\[
 \E_{\mathcal F}\,\pi(V_2)
 =\PP_{\mathcal F\otimes\pi}(x\in V_2)\le48N^2q_n .
\]

\paragraph{Step 4: choice of one realization.}
Let $V=V_1\cup V_2$. Combining \eqref{eq:ref-miss-mass} with
the two bounds of Step~3, we conclude that
\[
 \E_{\mathcal F}\bigl[\mathbf 1_{\mathcal O}+\pi(V)\bigr]
 \le2e^{-4L}+\bigl(8N^3+48N^2\bigr)q_n
 \le2e^{-4L}+56e^{3C_NL}q_n\le3e^{-4L},
\]
where the last inequality holds for all sufficiently large $n$ by
\eqref{eq:ref-absorb} with $C=3C_N+4$. On the finite probability space of
Step~1, some realization of the lists satisfies
$\mathbf 1_{\mathcal O}+\pi(V)\le\E_{\mathcal F}[\mathbf 1_{\mathcal O}+\pi(V)]\le3e^{-4L}$.
We fix it. Then $\pi(V)\le3e^{-4L}$, which is \eqref{b:dyn:fixed-bank}, and
the realization does not belong to $\mathcal O$, since
$\mathbf 1_{\mathcal O}\in\{0,1\}$ and $3e^{-4L}<1$.

\paragraph{Step 5: the four properties.}
Property~(i) holds because $V\supseteq V_1$. For~(ii), let $x$, $v$,
and ${\widetilde v}$ be as there. The triple $(v,x,{\widetilde v})$ is one of the orderings in the
definition of $V_2$, with $\sigma''=1$. Since $x\notin V_2$,
it does not lie in $\mathcal B_n\supseteq\mathcal B_n(q_*,b_1)$. By the
definition \eqref{eq:ref-fixed-violation} and the hypotheses of~(ii), this
means that $|R(v,x)-R(v,{\widetilde v})|\le2\eta_n$. The same argument with
$\mathcal N_n\subseteq\mathcal B_n$ proves~(iv). For~(iii), let $v$, $y$,
${\widetilde v}$, and $\sigma$ be as there. Then $R(y,\sigma {\widetilde v})=|R(y,{\widetilde v})|\ge b_1$,
$|R(v,y)|\le b_2$, and $|R(v,\sigma {\widetilde v})|\le b_2$. Since the realization does not belong to $\mathcal O$, the triple
$(v,y,\sigma {\widetilde v})$ of signed entries from three distinct levels does not lie in
$\mathcal B_n(b_1,b_2)$. This implies that
$|R(v,y)-\sigma R(v,{\widetilde v})|\le2\eta_n$, and the last claim of~(iii) follows from
$\bigl||R(v,y)|-|R(v,{\widetilde v})|\bigr|\le|R(v,y)-\sigma R(v,{\widetilde v})|$.

\paragraph{Step 6: the time scale $L_n=An^\lambda$.}
If $L_n=An^\lambda$ with $0<\lambda<1$, then $L_n\to\infty$ and
$L_n/n\to0$. For a fixed $C>0$, we have
\[
 e^{CL_n}q_n=C'(n+1)^3\Bigl(e^{CAn^\lambda-cR^3n^{1-3\lambda}}
   +e^{CAn^\lambda-cn^{1-2\xi}}+e^{CAn^\lambda-cn}\Bigr),
\]
with $C'$ the constant of \eqref{b:dyn:uniform-gibbs}. By
\eqref{eq:ref-exponents}, we have $1-3\lambda>\lambda$ and
$1-2\xi>1/2>\lambda$. Hence there exists $\lambda'>0$ such that each of the
three exponents is at most $-n^{\lambda'}$ for all large $n$, and
$(n+1)^3e^{-n^{\lambda'}}\to0$. This proves \eqref{eq:ref-absorb} and
completes the proof.
\end{proof}

The proof gives slightly more than (ii)--(iv). For every $x\notin V$, no
ordering of $x$ or $-x$ with two signed entries from distinct levels lies in
$\mathcal B_n$, and no ordering of three signed entries from distinct levels
lies in $\mathcal B_n$. Properties (ii)--(iv) are the instances of these
statements used in \cref{sec:induction}.

\begin{remark}[The same family for every path avoiding $V$]
\label{rem:ref-every-path}
The family $\mathcal F$ and the set $V$ of \cref{lem:ref-families} depend
only on $J$, and properties (i)--(iv) are deterministic statements about
$\mathcal F$, $V$, and configurations outside $V$. No path enters their
construction. We may therefore apply them along any finite sequence
$x_0,\ldots,x_m$ in $\Sigma_n\setminus V$ and for any choice of the
thresholds of \cref{sec:induction}: properties (i), (ii), and (iv) hold at
every $x_k$, simultaneously for all $k$, all levels, and all entries, and
(iii) holds for all entries. This is the form in which we use the lemma in
the induction of \cref{sec:induction}.

The order of choices matters for the stationary sign test of
\cref{sec:barrier-proof}. There, we bound the probability that the
stationary discrete heat-bath chain visits $V$ at one of the indices
$0,\ldots,m$ by $(m+1)\pi(V)$, using only that its law is $\pi$ at each fixed
integer index. This union bound requires that $V$ is fixed before the chain
is run. Note also that a
reference selected along a path depends on the path and need not have law
$\pi$; for this reason we state (ii)--(iv) for all entries of all lists, and
not only for selected references. Finally, we choose the family separately
at each $J\in\mathcal G_n$ and assert no measurable dependence on $J$.
None is needed, because the quantities that we integrate over the disorder
in \cref{sec:barrier-proof} do not depend on the family: they are the Gibbs
masses of initial states whose transition law is close to $\pi$ in total
variation.
\end{remark}
